\section*{Overview}

Binary search on answer is the pattern for optimization problems where the answer itself is not easy to construct
directly, but feasibility at a fixed candidate value \(\lambda\) \emph{is} easy to test.

\section*{When to Use It}

Use it when:

\begin{itemize}
  \item the statement asks for a minimum feasible value or maximum feasible value,
  \item there is a checker \(\texttt{can}(\lambda)\),
  \item the checker is monotone over the searched answer space.
\end{itemize}

The usual contest phrasing is ``minimize the maximum'', ``maximize the minimum'', or ``find the first time / capacity /
threshold where the plan becomes possible.''

\section*{Core Idea}

Turn the optimization target into a yes/no predicate:
\[
\texttt{can}(\lambda) = \text{``the problem is solvable if the answer is constrained by } \lambda \text{''}.
\]

Then prove that once the predicate becomes true, it stays true forever, or once it becomes false, it stays false
forever. Ordinary binary search finds the boundary.

\section*{Key Insight}

The search loop is almost never the real difficulty. The real work is modeling. A good note to self is:

\begin{quote}
Do not start writing the binary search until the predicate and its monotonicity are both explicit.
\end{quote}

\section*{Worked Problem}

\paragraph{Problem.}
An array of positive task durations must be partitioned into at most \(k\) consecutive days. Minimize the maximum total
duration assigned to one day.

\paragraph{Why binary search fits.}
Fix a candidate limit \(L\). The checker asks: can the tasks be split into at most \(k\) consecutive groups such that
every group sum is at most \(L\)?

\paragraph{Checker.}
Scan left to right and greedily extend the current day while the sum stays \(\le L\). When adding the next task would
exceed \(L\), start a new day.

\paragraph{Why the checker is correct.}
For positive tasks, postponing a split cannot hurt: every extra task only increases the current day's load. So the
greedy scan uses the minimum possible number of groups for this limit. If even that greedy grouping needs more than
\(k\) days, no other grouping can do better.

\paragraph{Why the predicate is monotone.}
If limit \(L\) is feasible, then every larger limit is also feasible because the same partition still works.

\section*{Problem Pattern}

This modeling trick reappears in:

\begin{itemize}
  \item maximum minimum distance under greedy placement,
  \item machine-capacity or scheduling thresholds,
  \item shortest time until a process can finish,
  \item DP or prefix-sum checkers that become feasible past one boundary.
\end{itemize}

\section*{Correctness Intuition}

Binary search itself contributes no special proof once the predicate is monotone. The proof burden sits entirely inside:

\begin{itemize}
  \item the checker correctness,
  \item the monotonicity argument,
  \item the bounds that guarantee the answer lies in the searched interval.
\end{itemize}

\section*{Complexity Analysis}

If the checker costs \(T\) and the answer space size is \(A\), the total cost is \(O(T \log A)\).

For the partition example, the checker is linear, so the full solution is \(O(n \log \sum a_i)\).

\section*{Implementation}

The code keeps:

\begin{itemize}
  \item a generic \texttt{first\_true} binary search,
  \item a concrete checker for the partition problem,
  \item a helper that returns the minimum feasible limit.
\end{itemize}

\section*{Common Pitfalls}

\begin{itemize}
  \item Searching before proving monotonicity.
  \item Picking bounds that do not actually contain the answer.
  \item Using a slow checker so the outer \(\log\) factor is irrelevant compared to a better direct solution.
  \item Applying integer binary search to a real-valued optimization problem without a precision plan.
\end{itemize}

\section*{Variants / Extensions}

\begin{itemize}
  \item Real-valued answer search.
  \item Parallel binary search for many offline queries.
  \item Parametric search where the checker itself is another greedy or DP routine.
\end{itemize}

\section*{Practice Problems}

\begin{itemize}
  \item Partition an array while minimizing the largest segment sum.
  \item Place routers or cows while maximizing the minimum distance.
  \item Find the smallest machine speed, budget, or threshold that makes a plan feasible.
\end{itemize}

\section*{References}

\begin{itemize}
  \item \href{https://cp-algorithms.com/num_methods/binary_search.html}{cp-algorithms: Binary Search on the Answer}.
  \item \href{https://usaco.guide/silver/binary-search}{USACO Guide: Binary Search}.
  \item \href{https://codeforces.com/catalog}{Codeforces Catalog}.
\end{itemize}
